\chapter{Conclusion and Open Research}
\section{Established contribution}
SPARK has an inspectable offline CNN/quantization pipeline, concrete host gate and event-encoding logic, and gateway storage/reporting/API software. Archived August and October runs provide preliminary fall-classification results with explicit operating-point tradeoffs. Their distinct splits, model artifacts, and thresholds prevent presenting them as one final validated detector.

The existing two-second, 100 Hz input contract resolves the report-level sampling contradiction. It does not supply the physical count/units adapter or continuous acquisition needed by the intended 200 Hz logger. The newer model bytes are available in the firmware header, while real TFLite Micro invocation remains unfinished.

\section{Limits of the present draft}
No completed wrist-worn cohort, elderly clinical study, measured battery endurance, physical BLE evaluation, or integrated on-device detection result is established here. Gradient-times-input and peak-feature explanations have different inputs and meanings; neither should be promoted to exact SHAP or causal clinical evidence. The available specificity tradeoffs and quantization degradation must remain visible during threshold selection.

The final thesis cannot declare the human system complete until acquisition, approval, collection, and integrated evaluation gates are met. Failing to meet a performance target is an experimental outcome to analyze, not a reason to rename a threshold or merge unrelated metrics.

\section{Next experiments}
First recover immutable split and model provenance and resolve the October threshold disagreement. Then validate the logger and real interpreter after procurement and parts receipt. Collect only under institutional approval, evaluate all annotated falls and non-fall exposure, and measure device timing/power directly. These experiments determine whether the proposed design achieves its human objective.

The separate robotics question is whether calibrated forearm signals support held-out monitoring after domain adaptation. Stable motion and independent annotations are prerequisites. Simulation results would still leave physical robot validation and mitigation-control safety unresolved.
